\documentclass[12pt,a4paper]{article}

\usepackage[T2A]{fontenc}
\usepackage[utf8]{inputenc}
\usepackage[russian]{babel}
\usepackage{amsmath,amssymb,amsthm}
\usepackage{array}
\usepackage{booktabs}
\usepackage{enumitem}
\usepackage{tikz}
\usepackage{listings}
\usepackage{xcolor}
\usepackage{geometry}

\geometry{margin=2.2cm}

\newtheorem{definition}{Определение}
\newtheorem{theorem}{Теорема}
\newtheorem{remark}{Замечание}

\lstset{
    basicstyle=\ttfamily\small,
    frame=single,
    breaklines=true,
    columns=fullflexible,
    keywordstyle=\color{blue},
    commentstyle=\color{gray}
}

\title{Билет №68 \\ \small Одноуровневые и многоуровневые дисциплины циклического обслуживания процессов на центральном процессоре, выбор кванта. (ИТМО)}
\author{Сериков Владислав Александрович}
\date{16 мая 2026}

\begin{document}

\maketitle
\tableofcontents
\newpage

\section{Базовые понятия планирования процессора}

\begin{definition}
\textbf{Процесс} --- выполняющаяся программа вместе со своим адресным пространством, состоянием регистров, открытыми файлами и другой системной информацией.
\end{definition}

\begin{definition}
\textbf{Планирование процессора} --- выбор одного из готовых к выполнению процессов или потоков для передачи ему центрального процессора.
\end{definition}

В однопроцессорной системе в каждый момент времени реально исполняется не более одного процесса. Если процессов несколько, операционная система создает иллюзию параллельного выполнения за счет быстрого переключения процессора между ними.

Основные состояния процесса:
\[
\text{new} \rightarrow \text{ready} \rightarrow \text{running} \rightarrow \text{terminated}.
\]

Также процесс может переходить из состояния \texttt{running} в состояние \texttt{waiting}, если ожидает завершения операции ввода-вывода, события, освобождения ресурса и т.п.

\begin{definition}
\textbf{Очередь готовых процессов} --- множество процессов, находящихся в состоянии \texttt{ready} и ожидающих выделения процессора.
\end{definition}

\begin{definition}
\textbf{Диспетчер} --- компонент операционной системы, непосредственно передающий процессор выбранному процессу: выполняет переключение контекста, переход в пользовательский режим, переход к нужной инструкции процесса.
\end{definition}

\begin{definition}
\textbf{Переключение контекста} --- сохранение состояния текущего процесса и восстановление состояния другого процесса.
\end{definition}

Переключение контекста не является полезной работой с точки зрения пользовательских процессов. Поэтому слишком частые переключения снижают производительность системы.

\section{Критерии качества планирования}

Для оценки дисциплин планирования используют следующие показатели.

\begin{enumerate}
    \item \textbf{Загрузка процессора}:
    \[
    U = \frac{T_{\text{полезная работа}}}{T_{\text{наблюдения}}}.
    \]

    \item \textbf{Пропускная способность} --- число процессов, завершенных за единицу времени:
    \[
    X = \frac{N}{T}.
    \]

    \item \textbf{Время оборота} процесса:
    \[
    T_{\text{turnaround}} = T_{\text{finish}} - T_{\text{arrival}}.
    \]

    \item \textbf{Время ожидания}:
    \[
    T_{\text{waiting}} = T_{\text{turnaround}} - T_{\text{CPU}},
    \]
    где $T_{\text{CPU}}$ --- суммарное время фактического выполнения процесса на процессоре.

    \item \textbf{Время отклика}:
    \[
    T_{\text{response}} = T_{\text{first\_run}} - T_{\text{arrival}}.
    \]

    \item \textbf{Справедливость} --- отсутствие неоправданного предпочтения одних процессов другим.

    \item \textbf{Отсутствие голодания} --- гарантия, что каждый готовый процесс рано или поздно получит процессор.
\end{enumerate}

Для интерактивных систем особенно важно малое время отклика, а для пакетных систем --- высокая пропускная способность и малое среднее время оборота.

\section{Вытесняющее и невытесняющее планирование}

\begin{definition}
\textbf{Невытесняющее планирование} --- дисциплина, при которой процесс, получивший процессор, выполняется до добровольной уступки процессора: завершения, блокировки на ввод-вывод или явного системного вызова.
\end{definition}

\begin{definition}
\textbf{Вытесняющее планирование} --- дисциплина, при которой операционная система может принудительно отнять процессор у выполняющегося процесса.
\end{definition}

Циклическое обслуживание процессов относится к вытесняющим дисциплинам. Его основная идея состоит в том, что каждый процесс получает процессор не навсегда, а на ограниченный интервал времени.

\section{Квант времени}

\begin{definition}
\textbf{Квант времени} $q$ --- максимальный непрерывный интервал времени, в течение которого процесс может выполняться на процессоре до принудительного вытеснения.
\end{definition}

Если процесс завершился или заблокировался раньше истечения кванта, остаток кванта не передается другому процессу. Если процесс не завершился за квант, он вытесняется и помещается обратно в очередь готовых процессов.

Механизм реализации кванта обычно связан с аппаратным таймером. Операционная система устанавливает таймер на величину $q$. После истечения кванта возникает прерывание таймера, управление переходит в ядро, и планировщик выбирает следующий процесс.

\section{Одноуровневое циклическое обслуживание}

\subsection{Дисциплина Round Robin}

\begin{definition}
\textbf{Round Robin} --- одноуровневая циклическая дисциплина планирования, при которой все готовые процессы находятся в одной очереди FIFO, а каждый процесс получает процессор не более чем на один квант времени $q$.
\end{definition}

Идея Round Robin:
\begin{itemize}
    \item все процессы равноправны;
    \item очередь готовых процессов обслуживается циклически;
    \item после истечения кванта незавершенный процесс становится в конец очереди;
    \item если процесс завершился или заблокировался раньше, планировщик сразу выбирает следующий процесс.
\end{itemize}

\subsection{Алгоритм Round Robin}

Пусть имеется очередь готовых процессов $Q$.

\begin{enumerate}
    \item Если очередь $Q$ пуста, процессор простаивает или выполняется процесс ядра.
    \item Извлечь из головы очереди первый процесс $P$.
    \item Выделить процессу $P$ процессор на время
    \[
    \min(q, r_P),
    \]
    где $r_P$ --- оставшееся процессорное время процесса.
    \item Если процесс $P$ завершился, удалить его из системы.
    \item Если процесс $P$ заблокировался до истечения кванта, перевести его в состояние ожидания.
    \item Если квант истек, а процесс $P$ не завершился и не заблокировался, вытеснить его и поместить в конец очереди $Q$.
    \item Перейти к шагу 1.
\end{enumerate}

Псевдокод:

\begin{lstlisting}
while true:
    if ready_queue.empty():
        idle()
    else:
        P = ready_queue.pop_front()
        run P for at most q time units

        if P.finished():
            destroy P
        else if P.blocked():
            move P to waiting_queue
        else:
            ready_queue.push_back(P)
\end{lstlisting}

\subsection{Минимальный пример Round Robin}

Пусть все процессы поступили в момент времени $0$, квант $q=2$.

\[
\begin{array}{c|c}
\toprule
\text{Процесс} & \text{Требуемое CPU-время} \\
\midrule
P_1 & 5 \\
P_2 & 3 \\
P_3 & 1 \\
\bottomrule
\end{array}
\]

Очередь в начале:
\[
Q = [P_1, P_2, P_3].
\]

Ход выполнения:

\[
\begin{array}{c|c|c}
\toprule
\text{Интервал времени} & \text{Выполняется процесс} & \text{Остатки после интервала} \\
\midrule
0-2 & P_1 & P_1:3,\; P_2:3,\; P_3:1 \\
2-4 & P_2 & P_1:3,\; P_2:1,\; P_3:1 \\
4-5 & P_3 & P_1:3,\; P_2:1,\; P_3:0 \\
5-7 & P_1 & P_1:1,\; P_2:1 \\
7-8 & P_2 & P_2:0,\; P_1:1 \\
8-9 & P_1 & P_1:0 \\
\bottomrule
\end{array}
\]

Диаграмма Ганта:

\[
\begin{array}{c|c|c|c|c|c|c}
0 & 2 & 4 & 5 & 7 & 8 & 9 \\
\hline
P_1 & P_2 & P_3 & P_1 & P_2 & P_1
\end{array}
\]

Более наглядно:

\[
\begin{tikzpicture}
\draw (0,0) rectangle (2,0.8);
\draw (2,0) rectangle (4,0.8);
\draw (4,0) rectangle (5,0.8);
\draw (5,0) rectangle (7,0.8);
\draw (7,0) rectangle (8,0.8);
\draw (8,0) rectangle (9,0.8);

\node at (1,0.4) {$P_1$};
\node at (3,0.4) {$P_2$};
\node at (4.5,0.4) {$P_3$};
\node at (6,0.4) {$P_1$};
\node at (7.5,0.4) {$P_2$};
\node at (8.5,0.4) {$P_1$};

\node at (0,-0.3) {$0$};
\node at (2,-0.3) {$2$};
\node at (4,-0.3) {$4$};
\node at (5,-0.3) {$5$};
\node at (7,-0.3) {$7$};
\node at (8,-0.3) {$8$};
\node at (9,-0.3) {$9$};
\end{tikzpicture}
\]

Времена завершения:
\[
C_1 = 9,\quad C_2 = 8,\quad C_3 = 5.
\]

Так как все процессы поступили в момент $0$, времена оборота:
\[
T_1 = 9,\quad T_2 = 8,\quad T_3 = 5.
\]

Времена ожидания:
\[
W_1 = 9 - 5 = 4,
\]
\[
W_2 = 8 - 3 = 5,
\]
\[
W_3 = 5 - 1 = 4.
\]

Среднее время ожидания:
\[
\overline{W} = \frac{4+5+4}{3} = \frac{13}{3} \approx 4.33.
\]

\subsection{Свойства Round Robin}

\subsubsection{Справедливость}

Round Robin считается справедливой дисциплиной, поскольку при постоянном наборе готовых процессов каждый из них регулярно получает процессор.

Если в очереди $n$ готовых процессов и все они используют полный квант, то каждый процесс получает процессор один раз за цикл длиной приблизительно
\[
nq.
\]

С учетом накладных расходов на переключение контекста $s$ длительность полного цикла:
\[
n(q+s).
\]

\subsubsection{Ограниченность времени ожидания до следующего обслуживания}

\begin{theorem}
Пусть в одноуровневой системе Round Robin имеется $n$ готовых процессов, все они всегда готовы к выполнению, квант равен $q$, временем переключения контекста пренебрегаем. Тогда время между двумя последовательными выделениями процессора одному и тому же процессу не превосходит
\[
(n-1)q.
\]
\end{theorem}

\begin{proof}
Рассмотрим произвольный процесс $P$. После того как $P$ отработал свой квант, он помещается в конец очереди. Перед ним в очереди могут находиться не более $n-1$ других процессов. Каждый из этих процессов может выполняться не более одного кванта $q$, после чего либо завершится, либо будет помещен в конец очереди.

Следовательно, до следующего запуска процесса $P$ суммарно могут выполниться не более $n-1$ процессов, каждый не более $q$ времени. Поэтому максимальное время ожидания:
\[
q + q + \ldots + q = (n-1)q.
\]
Теорема доказана.
\end{proof}

Если учитывать время переключения контекста $s$, то оценка становится:
\[
(n-1)(q+s).
\]

\subsubsection{Отсутствие голодания}

\begin{theorem}
Если множество готовых процессов конечно, каждый процесс не может бесконечно долго удерживать процессор сверх кванта $q$, а планировщик работает по дисциплине Round Robin, то ни один постоянно готовый процесс не будет голодать.
\end{theorem}

\begin{proof}
Пусть в очереди готовых процессов находится конечное число процессов $n$. Рассмотрим постоянно готовый процесс $P$. После каждого своего выполнения он либо завершится, либо будет помещен в конец очереди. Перед ним находится конечное число процессов, не более $n-1$.

Каждый из процессов перед ним может занимать процессор не более $q$ времени до вытеснения. Значит, через конечное время, не превосходящее $(n-1)q$ без учета переключений контекста, процесс $P$ снова окажется в голове очереди и получит процессор.

Следовательно, невозможно, чтобы процесс $P$ оставался готовым бесконечно долго и никогда не получал процессор. Значит, голодание отсутствует.
\end{proof}

\section{Влияние размера кванта}

Размер кванта $q$ существенно влияет на поведение Round Robin.

\subsection{Слишком малый квант}

Если квант слишком мал, то переключения контекста происходят часто.

Пусть:
\[
q = \text{квант времени},
\]
\[
s = \text{время переключения контекста}.
\]

Если каждый процесс использует полный квант, то на один цикл обслуживания процесса тратится:
\[
q+s.
\]

Доля полезного времени:
\[
\eta = \frac{q}{q+s}.
\]

Доля накладных расходов:
\[
\delta = \frac{s}{q+s}.
\]

Например, если:
\[
q = 1 \text{ мс}, \quad s = 0.1 \text{ мс},
\]
то:
\[
\eta = \frac{1}{1+0.1} = \frac{1}{1.1} \approx 0.909.
\]

Около $9.1\%$ времени теряется на переключения.

Если же:
\[
q = 10 \text{ мс}, \quad s = 0.1 \text{ мс},
\]
то:
\[
\eta = \frac{10}{10.1} \approx 0.990.
\]

Накладные расходы составляют около $1\%$.

\subsection{Слишком большой квант}

Если квант слишком велик, Round Robin приближается к FCFS.

\begin{definition}
\textbf{FCFS} --- дисциплина First-Come, First-Served, при которой процессы обслуживаются в порядке поступления без вытеснения.
\end{definition}

Если
\[
q \geq \max_i B_i,
\]
где $B_i$ --- полная длительность CPU-burst процесса $P_i$, то каждый процесс завершает свой CPU-burst за один запуск. В этом случае Round Robin фактически совпадает с FCFS для данного набора процессов.

Недостаток большого кванта: интерактивный процесс может долго ждать, пока завершится длинный CPU-bound процесс.

\subsection{Компромисс при выборе кванта}

При выборе кванта нужно соблюдать баланс:
\[
s \ll q,
\]
чтобы накладные расходы были малы, но при этом $q$ не должен быть настолько большим, чтобы ухудшать время отклика.

Обычно стремятся, чтобы:
\[
\frac{s}{q+s} \leq \varepsilon,
\]
где $\varepsilon$ --- допустимая доля накладных расходов.

Решим относительно $q$:
\[
\frac{s}{q+s} \leq \varepsilon.
\]

Тогда:
\[
s \leq \varepsilon(q+s),
\]
\[
s \leq \varepsilon q + \varepsilon s,
\]
\[
s(1-\varepsilon) \leq \varepsilon q,
\]
\[
q \geq \frac{s(1-\varepsilon)}{\varepsilon}.
\]

Например, если время переключения контекста:
\[
s = 0.1 \text{ мс},
\]
а допустимые накладные расходы:
\[
\varepsilon = 0.01,
\]
то:
\[
q \geq \frac{0.1(1-0.01)}{0.01} = \frac{0.099}{0.01} = 9.9 \text{ мс}.
\]

Значит, квант порядка $10$ мс обеспечивает примерно не более $1\%$ потерь на переключения.

\section{Одноуровневые дисциплины циклического обслуживания}

Под одноуровневой дисциплиной понимается схема, в которой все готовые процессы находятся в одной очереди или в одном классе обслуживания. Различий по уровням приоритета нет.

\subsection{Классический Round Robin}

Классический вариант:
\[
Q = [P_1,P_2,\ldots,P_n].
\]

Все процессы получают одинаковый квант:
\[
q_1 = q_2 = \ldots = q_n = q.
\]

Преимущества:
\begin{itemize}
    \item простота реализации;
    \item справедливое распределение CPU;
    \item отсутствие голодания при конечном числе процессов;
    \item хорошее время отклика при умеренном кванте.
\end{itemize}

Недостатки:
\begin{itemize}
    \item не различает интерактивные и вычислительные процессы;
    \item не учитывает приоритеты;
    \item среднее время ожидания может быть хуже, чем у SJF/SRTF;
    \item чувствителен к выбору кванта.
\end{itemize}

\subsection{Взвешенный Round Robin}

\begin{definition}
\textbf{Взвешенный Round Robin} --- вариант циклического обслуживания, при котором процессам назначаются веса $w_i$, влияющие на долю процессорного времени.
\end{definition}

Один из вариантов реализации: процесс $P_i$ получает квант
\[
q_i = w_i q_0,
\]
где $q_0$ --- базовый квант.

Тогда доля процессорного времени процесса $P_i$ при постоянной готовности всех процессов приблизительно равна:
\[
\frac{w_i}{\sum_{j=1}^{n} w_j}.
\]

\subsubsection{Алгоритм взвешенного Round Robin}

\begin{enumerate}
    \item Каждому процессу $P_i$ назначить вес $w_i$.
    \item Задать базовый квант $q_0$.
    \item Взять очередной процесс $P_i$ из очереди.
    \item Выделить ему процессор на время
    \[
    q_i = w_i q_0.
    \]
    \item Если процесс не завершился и не заблокировался, вернуть его в конец очереди.
\end{enumerate}

\subsubsection{Минимальный пример}

Пусть:
\[
q_0 = 2,\quad w_1 = 1,\quad w_2 = 2.
\]

Тогда:
\[
q_1 = 2,\quad q_2 = 4.
\]

Если оба процесса всегда готовы, цикл обслуживания:
\[
P_1(2),\; P_2(4),\; P_1(2),\; P_2(4),\ldots
\]

Доля процессора:
\[
P_1: \frac{1}{1+2} = \frac{1}{3},
\]
\[
P_2: \frac{2}{1+2} = \frac{2}{3}.
\]

\section{Многоуровневые дисциплины обслуживания}

Одноуровневый Round Robin плохо различает процессы разных типов. В реальных системах желательно по-разному обслуживать:
\begin{itemize}
    \item интерактивные процессы;
    \item фоновые вычислительные процессы;
    \item системные процессы;
    \item процессы реального времени;
    \item пакетные задания.
\end{itemize}

Для этого используются многоуровневые очереди.

\section{Многоуровневая очередь без обратной связи}

\begin{definition}
\textbf{Многоуровневая очередь} --- дисциплина планирования, в которой готовые процессы распределяются по нескольким очередям, каждая из которых имеет собственный приоритет и, возможно, собственный алгоритм обслуживания.
\end{definition}

Пусть имеется $k$ очередей:
\[
Q_0,Q_1,\ldots,Q_{k-1}.
\]

Обычно $Q_0$ имеет наивысший приоритет, а $Q_{k-1}$ --- низший.

Пример:
\[
Q_0: \text{системные процессы},
\]
\[
Q_1: \text{интерактивные процессы},
\]
\[
Q_2: \text{пакетные процессы}.
\]

\subsection{Статическое распределение процессов}

В многоуровневой очереди без обратной связи процесс при поступлении закрепляется за некоторой очередью и далее не меняет ее.

Например:
\[
P_i \in Q_j \quad \text{на все время существования процесса}.
\]

Внутри каждой очереди может использоваться своя дисциплина:
\[
Q_0: \text{Round Robin},
\]
\[
Q_1: \text{Round Robin},
\]
\[
Q_2: \text{FCFS}.
\]

\subsection{Выбор очереди для обслуживания}

Существует два основных способа выбора очереди.

\subsubsection{Фиксированные приоритеты между очередями}

Алгоритм:
\begin{enumerate}
    \item Найти непустую очередь с наибольшим приоритетом.
    \item Выбрать процесс из этой очереди по ее внутренней дисциплине.
    \item Если во время выполнения процесса появилась задача в более приоритетной очереди, возможна вытесняющая передача процессора более приоритетной задаче.
\end{enumerate}

Проблема: процессы в низкоприоритетных очередях могут голодать, если высокоприоритетные очереди постоянно непусты.

\subsubsection{Разделение времени между очередями}

Каждой очереди выделяется доля процессорного времени.

Например:
\[
Q_0: 60\%, \quad Q_1: 30\%, \quad Q_2: 10\%.
\]

Внутри выделенного времени каждая очередь обслуживает свои процессы по собственному алгоритму.

Такой подход уменьшает риск голодания низкоприоритетных очередей.

\section{Многоуровневая очередь с обратной связью}

\begin{definition}
\textbf{Многоуровневая очередь с обратной связью} \textbf{(Multilevel Feedback Queue, MLFQ)} --- дисциплина планирования, в которой процессы могут перемещаться между очередями в зависимости от своего поведения.
\end{definition}

Основная идея MLFQ:
\begin{itemize}
    \item интерактивные процессы, часто выполняющие ввод-вывод и быстро освобождающие CPU, должны иметь высокий приоритет;
    \item CPU-bound процессы, долго использующие процессор, постепенно опускаются в очереди с меньшим приоритетом;
    \item система не требует заранее знать длительность CPU-burst;
    \item приоритет процесса определяется динамически.
\end{itemize}

\subsection{Параметры MLFQ}

Многоуровневая очередь с обратной связью задается следующими параметрами:

\begin{enumerate}
    \item число очередей $k$;
    \item приоритет каждой очереди;
    \item алгоритм обслуживания внутри каждой очереди;
    \item квант времени для каждой очереди:
    \[
    q_0,q_1,\ldots,q_{k-1};
    \]
    \item правило помещения нового процесса в очередь;
    \item правило понижения приоритета;
    \item правило повышения приоритета;
    \item правило предотвращения голодания.
\end{enumerate}

\subsection{Типовая структура MLFQ}

Обычно:
\[
q_0 < q_1 < \ldots < q_{k-1}.
\]

Высокоприоритетные очереди имеют малые кванты, чтобы обеспечить малое время отклика. Низкоприоритетные очереди имеют большие кванты, чтобы уменьшить накладные расходы для длинных вычислительных процессов.

Пример:
\[
Q_0: q_0 = 8\text{ мс},
\]
\[
Q_1: q_1 = 16\text{ мс},
\]
\[
Q_2: q_2 = 32\text{ мс}.
\]

Или:
\[
Q_0: \text{Round Robin, } q=8,
\]
\[
Q_1: \text{Round Robin, } q=16,
\]
\[
Q_2: \text{FCFS}.
\]

\subsection{Базовый алгоритм MLFQ}

Пусть очереди упорядочены по убыванию приоритета:
\[
Q_0 > Q_1 > \ldots > Q_{k-1}.
\]

\begin{enumerate}
    \item Новый процесс помещается в очередь $Q_0$.
    \item Планировщик выбирает процесс из непустой очереди с наибольшим приоритетом.
    \item Если в одной очереди несколько процессов, они обслуживаются, например, по Round Robin.
    \item Если процесс использовал весь свой квант, он понижается в очередь меньшего приоритета:
    \[
    Q_i \rightarrow Q_{i+1}.
    \]
    \item Если процесс освободил процессор до истечения кванта, например из-за ввода-вывода, он остается на том же уровне или может быть повышен.
    \item Если процесс слишком долго ожидает, он повышается в более приоритетную очередь. Это называется старением.
\end{enumerate}

Псевдокод:

\begin{lstlisting}
on new process P:
    put P into Q[0]

scheduler:
    while true:
        i = highest_priority_non_empty_queue()
        if no such i:
            idle()
        else:
            P = Q[i].pop_front()
            run P for at most q[i]

            if P.finished():
                destroy P
            else if P.blocked_before_quantum_expired():
                keep P at same level or promote P
            else:
                if i < k - 1:
                    Q[i+1].push_back(P)
                else:
                    Q[i].push_back(P)
\end{lstlisting}

\subsection{Иллюстрация MLFQ}

Пусть есть три очереди:
\[
Q_0: q=2,
\]
\[
Q_1: q=4,
\]
\[
Q_2: q=8.
\]

Все новые процессы поступают в $Q_0$.

Пусть процесс $P$ требует $11$ единиц CPU-времени и не блокируется.

\[
\begin{array}{c|c|c|c}
\toprule
\text{Интервал} & \text{Очередь} & \text{Полученный квант} & \text{Остаток CPU-времени} \\
\midrule
0-2 & Q_0 & 2 & 9 \\
2-6 & Q_1 & 4 & 5 \\
6-11 & Q_2 & 5 \leq 8 & 0 \\
\bottomrule
\end{array}
\]

Процесс сначала получил высокий приоритет, но поскольку использовал весь квант, был понижен.

Диаграмма:

\[
\begin{tikzpicture}
\draw (0,0) rectangle (2,0.8);
\draw (2,0) rectangle (6,0.8);
\draw (6,0) rectangle (11,0.8);

\node at (1,0.4) {$P:Q_0$};
\node at (4,0.4) {$P:Q_1$};
\node at (8.5,0.4) {$P:Q_2$};

\node at (0,-0.3) {$0$};
\node at (2,-0.3) {$2$};
\node at (6,-0.3) {$6$};
\node at (11,-0.3) {$11$};
\end{tikzpicture}
\]

\subsection{Пример с интерактивным и вычислительным процессами}

Пусть:
\[
Q_0: q=2,\quad Q_1: q=4.
\]

Есть два процесса:
\[
P_1: \text{интерактивный, после 1 единицы CPU уходит на ввод-вывод},
\]
\[
P_2: \text{вычислительный, требует 10 единиц CPU}.
\]

Ход работы:
\begin{enumerate}
    \item Оба процесса поступают в $Q_0$.
    \item $P_1$ выполняется 1 единицу и блокируется, не исчерпав квант. Он сохраняет высокий приоритет.
    \item $P_2$ выполняется 2 единицы, исчерпывает квант и понижается в $Q_1$.
    \item После завершения ввода-вывода $P_1$ снова попадает в $Q_0$.
    \item Планировщик предпочитает $Q_0$, поэтому $P_1$ получает быстрый отклик.
    \item $P_2$ продолжает выполняться в $Q_1$, когда нет процессов в $Q_0$.
\end{enumerate}

Таким образом, MLFQ автоматически отличает интерактивные процессы от вычислительных.

\section{Старение и предотвращение голодания}

\begin{definition}
\textbf{Голодание} --- ситуация, при которой процесс находится в состоянии готовности, но неопределенно долго не получает процессор.
\end{definition}

В MLFQ голодание возможно, если в высокоприоритетных очередях постоянно появляются новые процессы. Тогда процессы в нижних очередях могут долго не выполняться.

\begin{definition}
\textbf{Старение} --- механизм постепенного повышения приоритета процесса, если он слишком долго ожидает обслуживания.
\end{definition}

Варианты старения:
\begin{enumerate}
    \item если процесс ждет более $T$ единиц времени, повысить его на один уровень;
    \item периодически переносить все процессы в верхнюю очередь;
    \item увеличивать динамический приоритет пропорционально времени ожидания.
\end{enumerate}

\subsection{Алгоритм MLFQ со старением}

\begin{enumerate}
    \item Все новые процессы помещаются в $Q_0$.
    \item Если процесс использовал весь квант на уровне $Q_i$, он переносится в $Q_{i+1}$, если $i < k-1$.
    \item Если процесс добровольно освободил CPU до истечения кванта, он остается на том же уровне.
    \item Если процесс ожидает более $T_{\text{age}}$, он повышается на один уровень:
    \[
    Q_i \rightarrow Q_{i-1}.
    \]
    \item Раз в $S$ единиц времени можно выполнить глобальное повышение:
    \[
    \forall P:\quad P \rightarrow Q_0.
    \]
\end{enumerate}

Псевдокод глобального повышения:

\begin{lstlisting}
every S milliseconds:
    for i in 1..k-1:
        move all processes from Q[i] to Q[0]
\end{lstlisting}

\begin{theorem}
Если в MLFQ используется периодическое глобальное повышение всех готовых процессов в верхнюю очередь каждые $S$ единиц времени, а очередь верхнего уровня обслуживается по Round Robin с конечным числом процессов и конечным квантом, то постоянно готовый процесс не будет голодать.
\end{theorem}

\begin{proof}
Рассмотрим произвольный постоянно готовый процесс $P$. Даже если он был понижен в самую низкую очередь, через время не более $S$ произойдет глобальное повышение, и $P$ будет помещен в верхнюю очередь $Q_0$.

После помещения в $Q_0$ очередь верхнего уровня обслуживается по Round Robin. При конечном числе процессов в $Q_0$ и конечном кванте каждый процесс из $Q_0$ получает процессор за конечное время. Это следует из свойства Round Robin: перед процессом находится конечное число процессов, каждый может занимать процессор не более одного кванта.

Следовательно, после каждого глобального повышения процесс $P$ получает процессор за конечное время. Поэтому он не может оставаться готовым бесконечно долго без обслуживания. Голодание отсутствует.
\end{proof}

\section{Выбор кванта времени}

Выбор кванта зависит от целей системы.

\subsection{Факторы, влияющие на квант}

\begin{enumerate}
    \item \textbf{Стоимость переключения контекста} $s$.

    Чем больше $s$, тем больше должен быть квант, иначе система будет тратить много времени на переключения.

    \item \textbf{Требования к времени отклика}.

    Для интерактивных систем квант должен быть достаточно малым.

    \item \textbf{Характер нагрузки}.

    Если много коротких интерактивных задач, полезны малые кванты на верхних уровнях. Если много CPU-bound задач, слишком малые кванты вредны.

    \item \textbf{Число одновременно готовых процессов} $n$.

    В Round Robin максимальное ожидание до следующего запуска примерно:
    \[
    (n-1)q.
    \]

    Поэтому при большом $n$ увеличение $q$ ухудшает время отклика.

    \item \textbf{Точность аппаратного таймера}.

    Квант не должен быть меньше разумно измеряемого и обслуживаемого интервала времени.

    \item \textbf{Кэш-поведение}.

    Частые переключения ухудшают локальность кэшей и TLB.
\end{enumerate}

\subsection{Нижняя граница кванта по накладным расходам}

Пусть допустимая доля накладных расходов:
\[
\varepsilon.
\]

Тогда:
\[
\frac{s}{q+s} \leq \varepsilon.
\]

Отсюда:
\[
q \geq \frac{s(1-\varepsilon)}{\varepsilon}.
\]

Это дает нижнюю оценку кванта.

\subsection{Верхняя граница кванта по времени отклика}

Если в системе $n$ готовых процессов, то грубая оценка времени ожидания до следующего обслуживания:
\[
T_{\text{wait}} \leq (n-1)q.
\]

Если требуется, чтобы задержка не превосходила $R_{\max}$, то:
\[
(n-1)q \leq R_{\max}.
\]

Следовательно:
\[
q \leq \frac{R_{\max}}{n-1}.
\]

Эта формула дает верхнюю оценку кванта.

\subsection{Совместный выбор кванта}

Нужно выбрать $q$, удовлетворяющий:
\[
\frac{s(1-\varepsilon)}{\varepsilon} \leq q \leq \frac{R_{\max}}{n-1}.
\]

Если такой $q$ существует, можно подобрать квант, одновременно обеспечивающий приемлемые накладные расходы и приемлемое время отклика.

Если интервал пуст:
\[
\frac{s(1-\varepsilon)}{\varepsilon} > \frac{R_{\max}}{n-1},
\]
значит, требования слишком жесткие для данной системы. Возможные решения:
\begin{itemize}
    \item уменьшить стоимость переключения контекста;
    \item уменьшить число одновременно активных процессов;
    \item ослабить требования к времени отклика;
    \item использовать многоуровневое планирование;
    \item выделить интерактивным задачам отдельные высокоприоритетные очереди.
\end{itemize}

\subsection{Практические эвристики}

На практике часто используют следующие правила:
\begin{enumerate}
    \item квант должен быть существенно больше времени переключения контекста:
    \[
    q \gg s;
    \]

    \item значительная доля CPU-burst должна завершаться за один квант;

    \item для интерактивных процессов лучше малые кванты;

    \item для вычислительных процессов допустимы большие кванты;

    \item в MLFQ кванты обычно возрастают при понижении приоритета:
    \[
    q_0 < q_1 < \ldots < q_{k-1}.
    \]
\end{enumerate}

\section{Сравнение одноуровневого RR и MLFQ}

\[
\begin{array}{p{4cm}|p{5cm}|p{5cm}}
\toprule
\textbf{Критерий} & \textbf{Round Robin} & \textbf{MLFQ} \\
\midrule
Число очередей & Одна & Несколько \\
\midrule
Приоритеты & Обычно отсутствуют & Динамические приоритеты \\
\midrule
Квант & Один общий квант & Разные кванты на разных уровнях \\
\midrule
Интерактивные процессы & Не выделяются специально & Обычно получают высокий приоритет \\
\midrule
CPU-bound процессы & Обслуживаются наравне со всеми & Постепенно понижаются \\
\midrule
Простота реализации & Высокая & Ниже \\
\midrule
Риск голодания & Обычно отсутствует при конечной очереди & Возможен без старения \\
\midrule
Гибкость & Низкая & Высокая \\
\bottomrule
\end{array}
\]

\section{Типичные ошибки при реализации циклического обслуживания}

\begin{enumerate}
    \item \textbf{Слишком малый квант.}

    Ведет к чрезмерным переключениям контекста.

    \item \textbf{Слишком большой квант.}

    Система ведет себя почти как FCFS, ухудшается интерактивность.

    \item \textbf{Отсутствие старения в MLFQ.}

    Низкоприоритетные процессы могут голодать.

    \item \textbf{Неправильный учет блокировок.}

    Если интерактивный процесс блокируется до истечения кванта, его не следует наказывать как CPU-bound процесс.

    \item \textbf{Нечестное восстановление остатка кванта.}

    Некоторые реализации учитывают, сколько процесс реально использовал CPU. Если это не делать аккуратно, возможны манипуляции со стороны процессов.

    \item \textbf{Игнорирование стоимости кэш-промахов.}

    Частые переключения влияют не только на время сохранения регистров, но и на состояние кэшей, TLB и конвейера процессора.
\end{enumerate}

\section{Итоговые положения}

\begin{enumerate}
    \item Циклическое обслуживание основано на выделении процессора процессам на ограниченный квант времени.

    \item Одноуровневый Round Robin использует одну очередь готовых процессов и одинаковый квант для всех.

    \item Round Robin прост, справедлив и обеспечивает отсутствие голодания при конечном числе готовых процессов.

    \item Слишком малый квант увеличивает накладные расходы на переключения контекста.

    \item Слишком большой квант превращает Round Robin в подобие FCFS и ухудшает время отклика.

    \item Многоуровневые очереди позволяют разделить процессы на классы.

    \item MLFQ динамически изменяет приоритет процесса в зависимости от его поведения.

    \item Интерактивные процессы в MLFQ обычно остаются на верхних уровнях, а CPU-bound процессы постепенно опускаются вниз.

    \item Для предотвращения голодания в MLFQ применяют старение и периодическое повышение приоритетов.

    \item Разумный выбор кванта должен учитывать стоимость переключения контекста, число готовых процессов, требования к времени отклика и характер нагрузки.
\end{enumerate}

\end{document}
